\subsection{Offline Preparation of the MANO Model}\label{subsec:model_prep}

The official \acrshort{mano} model is distributed as a legacy Python pickle that
depends on the \texttt{chumpy} autodifferentiation library and an old NumPy
object layout. Loading it directly on every run would be slow and would couple the
runtime to deprecated dependencies, so the model is converted once into a plain
NumPy dictionary that the pipeline loads thereafter.

\texttt{config} resolves the relevant paths relative to the package location. The
source directory holds the official releases (\texttt{MANO\_LEFT.pkl} and
\texttt{MANO\_RIGHT.pkl}); the output directory, created on first run if absent,
receives the converted files (\texttt{mano\_left.pkl} and \texttt{mano\_right.pkl}).

\texttt{prepare\_model} performs the conversion for both hands. It reads the
official pickle with \texttt{latin1} encoding and extracts only the fields the
runtime requires, renaming them to the project's schema: the pose
\acrshort{pca} basis and mean (from \texttt{hands\_components} and
\texttt{hands\_mean}), the joint regressor (densified from its sparse form via
\texttt{toarray}), the skinning weights, the pose and shape blend-shape bases
(from \texttt{posedirs} and \texttt{shapedirs}), the template mesh, the face
topology, and the kinematic tree (the first row of \texttt{kintree\_table}, with
the root's parent set to \texttt{None}). Because the original pickle was written
against an earlier NumPy and Python, the module first installs a small set of
compatibility aliases --- restoring removed type names such as \texttt{np.bool}
and \texttt{np.int}, and mapping \texttt{inspect.getargspec} to its modern
replacement --- so that the legacy object can be unpickled under the current
environment. The result is a self-contained file that loads quickly and carries
none of the original model's dependencies.
